% Source to translate: moto-touring-copilot/docs/proposal/source/04-proposed-scheme.md
% The architecture figure in the Chinese proposal has Chinese labels; redraw it in
% English and save as figures/architecture.pdf (the placeholder disappears once it exists).
\section{Proposed Scheme}
\label{sec:scheme}

\subsection{System Architecture}

The system follows an eight-module architecture (Fig.~\ref{fig:arch}).

\begin{figure*}[!t]
\centering
\figureorplaceholder{architecture.pdf}{width=0.9\textwidth}{6cm}
\caption{Architecture of Moto Touring Copilot (modules \ding{172}--\ding{179}).}
\label{fig:arch}
\end{figure*}

\begin{itemize}
  \item[\ding{172}] \textbf{User and scenario:} heavy-motorcycle touring riders in Taiwan planning a trip and checking weather, road conditions, fuel, and refueling points before departure.
  \item[\ding{173}] \textbf{Query input:} a web chat interface; the rider describes the trip by place names or picks points on a map, and imports GPX tracks to build and update the personal pace model.
  \item[\ding{174}] \textbf{Planner and reasoning:} five agents divide the work; a failed verification triggers replanning; rider preferences and past plans are remembered.
  \item[\ding{175}] \textbf{Data acquisition:} the rider's GPX tracks, CWA, TDX, the national digital terrain model, OpenStreetMap, fuel-station data, and the Google Maps Routes API.
  \item[\ding{176}] \textbf{Analytics engine:} the personal pace model (P50/P90), weather and road risk rules, fuel-range and station checks, and daylight checks, all implemented as deterministic Python tools.
  \item[\ding{177}] \textbf{Domain knowledge:} retrieval over CWA warning definitions, highway traffic information, riding-safety material, motorcycle specifications (tank size, fuel consumption), and related research, with sources cited.
  \item[\ding{178}] \textbf{Decision and recommendation:} recommended departure time, arrival-time ranges per segment, fuel stops, and ranked risks with reasons; infeasible requests are explained with alternatives.
  \item[\ding{179}] \textbf{Output:} a chat answer, a map dashboard with segments colored by time and risk, fuel stations, and sunset time, and a downloadable trip plan.
\end{itemize}

\subsection{Multi-Agent Design}
\todo{Orchestrator, pace, weather, road-and-supply, and verifier agents; verification rules 1--5 (source 4.4.1).}

\subsection{Personalized Pace Model}
\todo{Segment features, log moving-time target, quantile loss, conformal correction of trip-level P90 (sources 4.5--4.6).}

\subsection{Planning Algorithm}

\begin{algorithm}[!t]
\caption{\todo{Verification-driven planning}}
\label{alg:plan}
\begin{algorithmic}[1]
\Require \todo{route, date, rider GPX history}
\Ensure \todo{verified trip plan}
\State \todo{Translate the pseudocode in source 4.8 line by line.}
\end{algorithmic}
\end{algorithm}

\todo{Explain Algorithm~\ref{alg:plan}, including the replanning loop and exceptions (sources 4.7--4.8).}
